\section{Software Implementation}
\label{sec:experiments_software}

This section details the experimental evaluation of the proposed method.
\gls{dl} and \gls{rl} are strongly empirical fields, where experimental evaluation plays an important role in assessing how methods perform in practice.
Particular attention is therefore given to the implementation of \gls{car}.
It relies on established libraries that provide efficient implementations of its main components.
The code is made publicly available as open-source to support reproducible replication of the reported results.\footnote{\url{https://github.com/Maxalaar/Clustering_Agent_Latent_Space}}

The \gls{car} implementation is organized as a pipeline of five stages.
Each stage performs a specific operation and produces the artifacts required by the subsequent stages.
The stages are implemented as independent executables allowing each stage to be rerun independently.
The whole pipeline is driven by a single configuration file, which is the input of the process and gathers the parameters of every stage.
The \gls{car} pipeline is illustrated in Figure~\ref{fig:car_pipeline}, and Table~\ref{tab:software_libraries} summarizes the libraries used throughout the pipeline.
The remainder of this section describes the purpose and implementation of each stage.

\begin{figure}[t]
    \centering
    \resizebox{0.92\textwidth}{!}{
        \begin{tikzpicture}[
            font=\small,
            stage/.style={
                draw, rounded corners=0.15cm, thick,
                minimum width=6cm, minimum height=1.0cm,
                align=center, fill=black!3,
            },
            artifact/.style={align=left, font=\footnotesize\itshape, text=black!65},
            result/.style={align=left, font=\footnotesize, text=black!75},
            flow/.style={-{Latex[length=2.2mm]}, thick},
        ]
            \node[stage] (s1) {\textbf{1. Initial policy training}};
            \node[stage, below=1.5cm of s1] (s2) {\textbf{2. Trajectory dataset generation}};
            \node[stage, below=1.5cm of s2] (s3) {\textbf{3. Surrogate policy training}};
            \node[stage, below=2.0cm of s3] (s4) {\textbf{4. Surrogate policy evaluation}};
            \node[stage, below=1.5cm of s4] (s5) {\textbf{5. Latent space analysis}};

            \draw[flow] (s1) -- (s2);
            \draw[flow] (s2) -- (s3);
            \draw[flow] (s3) -- (s4);
            \draw[flow] (s3.south) -- ++(0,-0.55cm) -| ([xshift=-1.6cm]s5.north);

            \node[artifact] at ($(s1)!0.5!(s2) + (0.3cm,0)$) {initial policy $\OriginalPolicy$};
            \node[artifact] at ($(s2)!0.5!(s3) + (0.3cm,0)$) {trajectory dataset (Parquet)};
            \node[artifact, anchor=west] at ($(s3.south) + (0.3cm,-0.4cm)$) {$n$ surrogate policies $\SurrogatePolicy{1}, \dots, \SurrogatePolicy{n}$};

            % Stage 5 also reuses the trajectory dataset.
            \draw[flow, dashed] (s2.west) -- ++(-1.7cm,0) |- ([yshift=1.5mm]s5.west);
            \node[artifact, anchor=east] at ($(s5.west) + (-1.9cm,0.9cm)$) {trajectory\\ dataset};

            \node[result, anchor=west] at ($(s4.east) + (0.6cm,0)$) {$\rightarrow$~\BehaviorReconstructionName};
            \node[result, anchor=west, text width=4.6cm] at ($(s5.east) + (0.6cm,0)$) {$\rightarrow$~\ClusteringUniversalityName, \AbstractionScoreName, clusters};
        \end{tikzpicture}
    }
    \caption{
        The \gls{car} pipeline. Each stage is an independent executable; stages communicate only through artifacts written to disk (solid arrows). The dashed arrow indicates an artifact reused by a later stage.
    }
    \label{fig:car_pipeline}
\end{figure}


\FloatBarrier

\begin{table}[t]
\hspace*{\fill}
\resizebox{\textwidth}{!}{
\begin{tblr}{|Q[c,m,3.4cm]|Q[c,m,3.5cm]|Q[l,m,8.5cm]|}
    \hline
    \textbf{Pipeline Stage} & \textbf{Library} & \SetCell{c}\textbf{Description} \\
    \hline
    \hline
    \SetCell[r=3]{c} \textbf{Initial policy training} & RLlib \cite{moritz2018ray, liang2018rllib} & Distributed \gls{rl} library used to train the initial policy. \\
    \hline
    & PyTorch \cite{paszke2019pytorch} & Deep learning framework implementing the policy and critic networks. \\
    \hline
    & Gymnasium \cite{towers2024gymnasium} & Common interface used to plug in any \gls{rl} environment. \\
    \hline
    \hline
    \SetCell[r=2]{c} \textbf{Trajectory dataset generation} & RLlib \cite{moritz2018ray, liang2018rllib} & Distributed environment runners used to collect trajectories in parallel. \\
    \hline
    & Apache Arrow \cite{arrow} & Columnar file format (Parquet) used to store the trajectory datasets. \\
    \hline
    \hline
    \textbf{Surrogate policy training} & PyTorch Lightning \cite{falcon2019lightning} & Structured training loop used for the behavior cloning of the surrogate policies. \\
    \hline
    \hline
    \textbf{Surrogate policy evaluation} & RLlib \cite{moritz2018ray, liang2018rllib} & Wraps each surrogate policy as an \gls{rl} policy to evaluate it in the environment. \\
    \hline
    \hline
    \textbf{Latent space analysis} & cuML \cite{raschka2020cuml} & \gls{gpu}-accelerated clustering algorithms used to analyze the latent representations. \\
    \hline
\end{tblr}
}
\hspace*{\fill}
\caption{Main software libraries used at each stage of the \gls{car} pipeline.}
\label{tab:software_libraries}
\end{table}


\subsection{Initial policy training}

The first stage trains the initial policy using RLlib \cite{moritz2018ray, liang2018rllib}.
RLlib provides scalable implementations of \gls{rl} algorithms and supports distributed training.
The use of RLlib also allows \gls{car} to be applied to environments implemented using the Gymnasium interface \cite{towers2024gymnasium}.
By default, the initial policy is trained using \gls{ppo} \cite{schulman2017proximalpolicyoptimizationalgorithms}.
Other RLlib algorithms, such as \gls{appo} \cite{espeholt2018impala}, \gls{dqn} \cite{mnih2013playingatarideepreinforcement}, and DreamerV3 \cite{hafner2023masteringdiversedomainsworld}, can also be selected through the same configuration.

The policy and critic networks are implemented using PyTorch \cite{paszke2019pytorch}.
PyTorch is the standard \gls{dl} framework and is widely used for developing and training \gls{nn}.
It provides \gls{gpu} acceleration and supports a wide range of \gls{nn} architectures.
This makes it well suited for implementing the different policy and critic architectures used in the experiments.

\subsection{Trajectory dataset generation}

The second stage builds the dataset used to train the surrogate policies.
RLlib is used to collect trajectories in parallel using multiple environment runners.
For each step, the dataset stores the observation together with the parameters of the action distribution produced by the initial policy.

Trajectories are stored as Parquet files \cite{arrow}.
Its columnar structure allows efficient data access and supports concurrent processing.
This is useful here because the datasets contain around one million observation--action pairs and cannot be loaded entirely into memory.

\subsection{Surrogate policy training}

The third stage trains the surrogate policies using \gls{bc} on the trajectory dataset.
This provides a supervised learning approach for replicating the initial policy.
The training procedure is implemented using PyTorch Lightning \cite{falcon2019lightning}.
PyTorch Lightning provides a simplified interface for training \gls{nn} with PyTorch.
It handles common training operations such as data loading, optimization, checkpointing, and \gls{gpu} execution.

\subsection{Surrogate policy evaluation}

The fourth stage evaluates how faithfully the surrogate policies reproduce the behavior of the initial policy.
This stage therefore requires both the surrogate policies and the initial policy.
Each surrogate is wrapped as an RLlib policy and evaluated in the original environment.
The initial policy is evaluated under the same conditions to provide the reference return.
The mean cumulative reward of each surrogate over multiple episodes is then compared with that of the initial policy.
Evaluation episodes are collected using parallel environment runners.

This stage computes the \BehaviorReconstructionName.
It therefore verifies that the surrogate policies reproduce the behavior of the initial policy sufficiently well before their \glspl{lr} are analyzed.
If a surrogate does not meet the required level of behavior reconstruction, its training can be repeated.

\subsection{Latent space analysis}

The fifth stage performs the main analysis of \gls{car}.
The trajectory dataset and trained surrogate policies are used to generate the \glspl{lr}.
These representations form the point clouds that are analyzed using clustering algorithms.

The clustering step relies on the cuML library \cite{raschka2020cuml}, which provides \gls{gpu}-accelerated clustering algorithms.
This reduces the computational cost and allows point clouds containing tens of thousands of representations to be analyzed.

The resulting partitions are used to compute the \ClusteringUniversalityName and the \AbstractionScoreName.
The stage also produces the material required for the qualitative analysis, including the clustered point clouds and the cluster assignment of each observation.
